\section{Errores en el cálculo numérico}
Se utiliza para:
\begin{itemize}
    \item Saber que tan confiable es un rdo
    \item detectar mal condicionamiento
    \item garantizar una tolerancia aceptable
\end{itemize}

\subsection{Fuentes de Error}
\subsubsection{Error inherente}
Errores por cálculos previos, mediciones físicas, simplificaciones del modelo matemático.
Ej: $\theta''\ +\ g\ sen\left(\theta\right)\ = \ 0\ \implies\ \theta''\ +\ g\ \boldsymbol{\theta}\ =\ 0$
Esto es solo válido para theta acercandose a cero.
\subsubsection{Error de redondeo}
Por representación finita de un número en una computadora, se redondea y se pierde exactitud.
\subsubsection{Error de truncamiento o discretización}
Corte de procesos iterativos o sumatorias infinitas convergentes, el conocido hasta ahora.
Se da al particionar dominios continuos.
\subsection{Error absoluto y relativo}
Error absoluto: $\epsilon_a=a\tilde{~}-a$


Error relativo: $\epsilon_r=\frac{e_a}{a}$


Como generalmente no se conoce 'a', se trabaja con cotas de errores.


$\Delta\ a$: Cota de error absoluto. $|\epsilon_a|\le\Delta\ a$


$\frac{\Delta\ a}{a\tilde{~}}$ Cota de error relativo. $|\epsilon_r|\le \frac{\Delta\ a}{a\tilde{~}}$

El resultado junto a su error se expresan de la forma:
$a = a\tilde{~}+\Delta\ a$

\subsubsection{¿Cómo seleccionar $\Delta\ a$?}
Un criterio:
\begin{itemize}
    \item Se obtiene un dígito para $\Delta\ a$.
    \item Se ubica el dígito en el $a\tilde{~}$ afectado por la cota.
    \item Los dígitos que están a la derecha se desechan, pues no tiene sentido mantener la exactitud sobre dígitos de los que no estoy seguro que sean así.
    \item Se redondea de forma simétrica al más próximo.
\end{itemize}
El error relativo es la suma de los errores de las variables.
Los errores se propagan al realizar operaciones.
